\chapter*{Vorlesungsverzeichnis}
\addcontentsline{toc}{chapter}{B. Vorlesungsverzeichnis}
\chaptermark{Vorlesungsverzeichnis}

Hier ist eine Übersicht, wann welche Vorlesung stattfand, ob diese bereits im Latex-Dokument integiert wurde, und ob Beispiele und Erklärungen ergänzt wurden. Zudem wird vermerkt, ob die Vorlesung vollständikg als AnkiDeck aufzufinden ist.

\begin{table}[h!]
    \centering
    \caption{Verzeichnis der Vorlesungen und Übersciht des Arbeitsstandes. Folgende Abkürzungen wurden genutzt: T = Transskipt, E = Ergänzt, A = Ankideck.}
    \begin{tabular}{l | l C C C | l}
        \toprule
        \textbf{Vorlesung} & \textbf{Kapitel/Abschnitt} & \textbf{T} & \textbf{E} & \textbf{A} & \textbf{Datum} \\
        \midrule
        Vorlesung 1 & Fourierreihen -- Grundlagen & \surd & \times & \times & 15.04.26 \\
        Vorlesung 2 & Fourierentwicklung & \surd & \times & \times & 17.04.26 \\
        Vorlesung 3 & Fourierreihen -- L\textsuperscript{2}-Konvergenz & \surd & \times & \times & 22.04.26 \\

        Vorlesung 4 & Normierte und metrische Räume & \surd & \times & \times & 24.04.26 \\
        Vorlesung 5 & Cauchy-Folgen, Vollständigkeit, Topologie & \surd & \times & \times & 29.04.26 \\
        Vorlesung 6 & Rand, Inneres, Abschluss, Kompaktheit & \surd & \times & \times & 06.05.26 \\
        Vorlesung 7 & Kompaktheit, Skalarprodukte, Gram--Schmidt & \surd & \times & \times & 08.05.26 \\
        
        Vorlesung 8 & Orthogonalität in $\mathbb{R}^n$; Stetigkeit & \surd & \times & \times & 13.05.26 \\
        Vorlesung 9 & Lineare Abbildungen & \surd & \times & \times & 15.05.26 \\
        Vorlesung 10 & Stetige Funktionen auf Kompakten & \surd & \times & \times & 20.05.26 \\
        Vorlesung 11 & Konvergenz von Funktionenfolgen; Fixpunktsatz & \surd & \times & \times & 22.05.26 \\
        
        Vorlesung 12 & Partielle Ableitungen; Totale Ableitungen & \surd & \times & \times & 27.05.26 \\
        Vorlesung 13 & Differenzierbarkeit, Kettenregel, Mittelwertsatz & \surd & \times & \times & 29.05.26 \\

        Vorlesung 14 & MWS; Taylor-Entwicklung & \surd & \times & \times & 03.07.26 \\
        Vorlesung 15 & Extremwertaufgaben, SIF & \surd & \times & \times & 05.07.26 \\
        Vorlesung 16 & SIF, Extrema unter Nebenbedingungen & \surd & \times & \times & 10.07.26 \\
        Vorlesung 17 & NOB, UMF & \times & \times & \times & 12.07.26 \\
        Vorlesung 18 & DGL, Satz von Peano & \times & \times & \times & 17.07.26 \\
        Vorlesung 19 & Eindeutigkeit von AWA-Lösungen & \times & \times & \times & 19.06.2026 \\
        Vorlesung 20 & Globale Existenz und Stabilität & \times & \times & \times & 24.06.2026 \\
        Vorlesung 21 & Systeme lineare DGL  & \times & \times & \times & 26.06.2026 \\
        Vorlesung 22 &  & \times & \times & \times & 30.06.2026 \\
        % Vorlesung 23 &  & \times & \times & \times & \\
        % Vorlesung 24 &  & \times & \times & \times & \\
        % Vorlesung 25 &  & \times & \times & \times & \\
        % Vorlesung 26 &  & \times & \times & \times & \\
        \bottomrule
    \end{tabular}
\end{table}
